\documentclass[11pt,a4paper]{article}
\usepackage[margin=2.2cm,headheight=16pt]{geometry}
\usepackage[T1]{fontenc}
\usepackage{lmodern}
\usepackage{xcolor}
\usepackage{graphicx}
\usepackage{booktabs}
\usepackage{tabularx}
\usepackage{fancyhdr}
\usepackage[hidelinks,colorlinks=true,linkcolor=febblue,urlcolor=febblue]{hyperref}
\definecolor{febblue}{HTML}{1C4687}
\definecolor{febgold}{HTML}{FDB515}
\setlength{\parskip}{4pt}
\pagestyle{fancy}
\fancyhf{}
\fancyhead[L]{\color{febblue}\textbf{FEB Autonomous}}
\fancyhead[R]{\color{febblue}1:10 car budget}
\fancyfoot[C]{\thepage}
\renewcommand{\headrule}{\color{febgold}\hrule width\headwidth height 0.8pt}

\title{\color{febblue}\LARGE A 1:10 Autonomous Car for the Team\\[4pt]\large Budget breakdown under \$2{,}000}
\author{Formula Electric at Berkeley, Autonomous Team}
\date{September 2026}

\begin{document}
\maketitle

\section*{Summary}

We ask for \textbf{\$1{,}750} to build one 1:10-scale autonomous car and a lab track, with a
minimum configuration at \textbf{\$1{,}100} if the budget is tighter. The car is the
RoboRacer-standard Traxxas Slash with a 2D lidar, a small computer and a motor controller
that reports wheel speed. It is what new members learn real robotics on, what the perception
and control work is tested on between FSAE seasons, and the first place our code meets real
dynamics. The team's simulator already exists and is in use; this document is about the
hardware, area by area, with the costs, the choices inside each area, and the realities that
come with them.

\begin{tabularx}{\linewidth}{Xrr}
\toprule
Area & Recommended & Minimum \\
\midrule
Chassis and drive & 500 & 330 \\
Motor controller & 110 & 110 \\
Compute & 250 & 130 \\
Lidar & 320 & 100 \\
Camera & 0 & 0 \\
Power & 230 & 200 \\
Mounting and wiring & 80 & 60 \\
Track and cones & 45 & 15 \\
Spares and tools & 60 & 40 \\
\midrule
Subtotal & 1{,}595 & 985 \\
Contingency (10\,\%) & 160 & 100 \\
\textbf{Total} & \textbf{1{,}755} & \textbf{1{,}085} \\
\bottomrule
\end{tabularx}

\smallskip
\noindent USD, September 2026 list prices, rounded. Verify at order time; the lidar moves most.

\section{Where we are}

\begin{itemize}
  \item \textbf{Simulator.} The team has a working simulator of this exact car, a ROS~2
  devkit, a starter driving stack and a leaderboard. Members already write and race code in
  it. It needs nothing from this budget. Its limit is that nothing in it is real: no grip
  change with the floor, no servo backlash, no battery sag, no lidar seeing chair legs.
  \item \textbf{Owned hardware.} An Intel RealSense depth camera, a RoboSense M1 3D lidar,
  and a motorless Tamiya 1:10 touring chassis. The camera goes on this car; the other two are
  discussed in section~\ref{sec:setaside}.
  \item \textbf{Space.} A lab corridor is enough for a 1:10 track. The Alameda test track is
  for larger vehicles and is not needed here.
\end{itemize}

\section{The car, area by area}

\subsection{Chassis and drive}

\begin{tabularx}{\linewidth}{Xrl}
\toprule
Option & USD & \\
\midrule
Traxxas Slash 4x4 VXL (TRA68086-4), brushless motor and ESC included & 500 & recommended \\
Traxxas Slash 2WD VXL (TRA58076-4), brushless & 330 & minimum \\
\bottomrule
\end{tabularx}

\smallskip
The Slash is the chassis every RoboRacer build and drawing assumes, and the one the
simulator models. The 4x4 is the standard: it carries the sensor deck without lifting the
front wheels and tolerates novice driving. The 2WD saves \$170 and is fine for lidar-only
driving on a flat floor, but it spins the rear under a heavy deck and is not what the drawings
fit; buy it only if the total must come under \$1{,}200. Either comes ready to run with a
brushless motor; the stock speed controller is replaced by the item below, and the radio and
receiver stay as the manual-override link.

\subsection{Motor controller}

\begin{tabularx}{\linewidth}{Xrl}
\toprule
Option & USD & \\
\midrule
Flipsky mini FSESC 6.7 (VESC-compatible) & 110 & recommended \\
VESC 6 MkVI (reference part) & 260 & not needed at this budget \\
\bottomrule
\end{tabularx}

\smallskip
A VESC-type controller is not optional: it is what gives the software wheel speed
(odometry), current limits, and a clean serial interface. The reference VESC costs \$150 more
for better cooling and support; the Flipsky runs the same firmware and is what most student
builds use.

\subsection{Compute}

\begin{tabularx}{\linewidth}{Xrl}
\toprule
Option & USD & \\
\midrule
NVIDIA Jetson Orin Nano developer kit & 250 & recommended \\
Raspberry Pi 5, 8~GB, active cooler, SD card, PSU & 130 & minimum \\
\bottomrule
\end{tabularx}

\smallskip
The Pi runs ROS~2 and a lidar-only driver comfortably and is the cheaper way to a first lap.
It cannot run camera perception at useful rates, so the day the team wants cone detection
from the RealSense, it needs the Jetson. Buying the Jetson first avoids buying twice; buying
the Pi first gets a car on the floor for \$120 less and keeps the Jetson as the upgrade.

\subsection{Lidar}

\begin{tabularx}{\linewidth}{Xrl}
\toprule
Option & USD & \\
\midrule
Slamtec RPLidar A2M12, 12~m, 12~Hz, 0.25\textdegree & 320 & recommended \\
Slamtec RPLidar A1M8, 12~m, 5 to 10~Hz, 1\textdegree & 100 & minimum \\
Hokuyo UST-10LX, 10~m, 40~Hz, 0.25\textdegree & 1{,}500 & the reference; out of this budget \\
\bottomrule
\end{tabularx}

\smallskip
The lidar decides how fast the car can safely go. The reference Hokuyo scans at 40~Hz and is
what the simulator models; it is also most of a \$3{,}000 build. The A2M12 at 12~Hz is a
sensible car for 2 to 3~m/s and is what the budget allows. The A1M8 is a \$100 sensor that
proves the pipeline and nothing more: 8~kHz sample rate, a degree of angular resolution,
noisy beyond 6~m. Spinning lidars are also the part that breaks in a crash; mount it behind
the bumper line.

\subsection{Camera}

The team's RealSense goes on the car at no cost. With the lidar it gives the sensor pair for
cone detection: the camera classifies cones by colour, the lidar gives range and bearing. This
is the perception work that carries most directly to the FSAE car, and it needs the Jetson.

\subsection{Power}

\begin{tabularx}{\linewidth}{Xrl}
\toprule
Item & USD & \\
\midrule
Two Traxxas 3S 5000~mAh LiPo packs & 170 & one drives, one charges \\
Balance charger & 50 & \\
LiPo safety bag & 10 & \\
5~V / 5~A buck converter, XT60 splitter, inline fuse & 30 & feeds computer and lidar from the drive pack \\
\bottomrule
\end{tabularx}

\smallskip
Minimum: one pack (\$85 less). The RoboRacer powerboard (\$100) is a cleaner version of the
converter and splitter with a safe-shutdown line; not needed to start.

\subsection{Mounting and wiring}

Laser-cut deck plates from the published RoboRacer drawings, cut in a campus makerspace
(about \$30 of acrylic), standoffs and fasteners (\$15), a powered USB hub and cables (\$35).
Minimum \$60 with 3D-printed brackets instead of a cut deck.

\subsection{Track and cones}

\begin{tabularx}{\linewidth}{Xrl}
\toprule
Item & USD & \\
\midrule
Cardboard walls: flattened boxes cut into 20~cm strips, taped in loops & 0 & \\
Floor tape (start line, lane marks) & 15 & \\
FSAE-style cones, 3D-printed at 1:10, 100 pieces & 15 & about 10~g of filament each \\
Foam pipe insulation on the corners & 15 & \\
\bottomrule
\end{tabularx}

\smallskip
Cardboard is enough for the lab: the lidar sees it as well as any wall, a 20~cm strip is
well above the scan plane, and a wall that folds when hit is kinder to the car. It has to be
re-taped each session and does not store. The cones are the part worth printing on day one,
because the cone track is where the camera work happens. \emph{Later, if the track is in
daily use, 30~m of 33~cm flexible aluminium duct (about \$130) is the permanent version.}

\subsection{Spares and tools}

Tyres (\$30), pinion and bumper (\$20), zip ties, tape, Velcro (\$10). The team's own tools
otherwise. A wall hit at 3~m/s breaks a bumper or a servo saver, not the car; a lidar hit
breaks the lidar.

\section{Options considered and set aside}
\label{sec:setaside}

\begin{itemize}
  \item \textbf{The team's Tamiya chassis.} A motorless touring-car kit. It would save the
  \$500 chassis line and cost a motor, controller and servo (\$200) plus a custom sensor deck
  with no drawings to copy. Its wheelbase, track width and ground clearance differ from the
  car the simulator models and from every published build, so nothing transfers and every
  mounting problem is solved alone. Net saving about \$300 for a semester of mechanical work
  that teaches no autonomy. It stays a bench mule for wiring and sensor practice.
  \item \textbf{A 1:4 or 1:5 car with the team's M1 lidar.} No maker sells a 1:4 road car;
  the nearest are 1:5 trucks. With the owned M1 and the Alameda track counted, a 1:5 build
  costs about \$3{,}300 to \$3{,}900: chassis \$1{,}300 to \$1{,}900, 8S batteries and
  charger \$500, a high-current controller \$350, a Jetson Orin NX class computer for the
  M1's point clouds \$700, deck, roll protection and kill switch \$450. Twice this budget,
  and it cannot be run on campus: a 10~kg car at 50~km/h needs a closed outdoor area, so
  every test is a session in Alameda rather than an afternoon in the lab, and it is not a
  car new members can learn on. The M1 itself (about 0.7~kg, a half-metre blind zone,
  120\textdegree{} forward view) is a full-size-vehicle sensor that does not fit a 1:10 car's
  mass or corridor. Both wait for a later budget, after the 1:10 programme has run.
\end{itemize}

\section{What the money buys, in order}

\begin{tabularx}{\linewidth}{lX}
\toprule
When & Milestone \\
\midrule
Order + 2 weekends & Car assembled from the RoboRacer build guide, radio drive, all sensors on ROS~2 \\
+ 3 weeks & Starter stack drives the cardboard track; simulator-to-car differences written down \\
+ 6 weeks & Members' own stacks on the car; internal time-attack; first cone-detection results from the camera \\
Semester end & Perception and control lessons carried into the FSAE autonomous work \\
\bottomrule
\end{tabularx}

\section{The realities}

\begin{itemize}
  \item \textbf{One car.} Access has to be scheduled and earned; the simulator's
  qualification rule is the gate. Expect queueing in the weeks before an internal event.
  \item \textbf{It will be crashed.} A spinning lidar on the nose of a car driven by student
  code is a matter of when. The contingency and the spares line are for this. A dead A2M12 is
  \$320; a dead Hokuyo would be \$1{,}500, which is the other reason it is not in this budget.
  \item \textbf{The cheap parts cap performance.} At 12~Hz the lidar and the Pi limit safe
  speed to about 2 to 3~m/s; the simulator car goes faster. Members will feel the gap; that is
  the point, but it should be said up front.
  \item \textbf{Time is the scarce resource, not money.} Two weekends of assembly, a month
  to a first clean lap, and someone has to own the car, the batteries and the charging
  procedure. If nobody owns it, it becomes a shelf ornament by week six.
  \item \textbf{Batteries need a procedure.} LiPo bag, supervised charging, no charging
  unattended, a written checklist before the first charge.
  \item \textbf{The floor matters.} Lab tile and carpet grip differently; the car will be
  tuned for one floor. A cardboard track moves between sessions; lap times will not be
  comparable week to week until the duct arrives.
  \item \textbf{Cardboard walls have a life of weeks.} Budget half an hour per session for
  the track.
  \item \textbf{Out of scope for this budget.} The 1:5 car with the M1, a second car, and
  physical competitions. All are possible later; none are needed to get value from this car.
\end{itemize}

\bigskip
\noindent RoboRacer build guide and drawings: \url{https://f1tenth.readthedocs.io}\\
Team platform (simulator, devkit, manuals): \url{https://github.com/Pranman1/feb-racing}

\end{document}
